% varanc-presence-mixdom.tex — Reduced-tuple variational MixDom
% ancestral state inference (presence + (domain, fragchar) only).
% To be \input from tkf.tex.

\section{Mixture-of-trees variational MixDom ancestral inference}
\label{sec:varanc-presence-mixdom}

This appendix generalises the variational ancestral-presence
reconstruction of Section~\ref{sec:varanc-presence} from the
TKF92-conditional-WFST-as-MaxEnt-GGI approximation to the labelled
MixDom model, with two further simplifications relative to a fully
labelled treatment:
\begin{itemize}[nosep]
\item the per-character substitution class $c$ is integrated out at
the model level via the per-column prior $\classdist_{f, c}$
(rather than carried as a variational latent);
\item the fragchar-boundary indicator $g$ and domain-end indicator
$e$ are absorbed into a \emph{reduced} per-character WFST kernel that
sums over the implicit fragment/domain bookkeeping at each step
(rather than tracked explicitly in the variational state).
\end{itemize}
The variational state per (internal node, MSA column) is therefore
$\{N, D\} \cup \{(d, f)\}$, $|\mathcal{Z}| = N_{\text{dom}} N_{\text{fr}} + 2$
— a $(4 N_{\text{cl}})\times$ reduction over the fully-labelled state
space. Three concrete payoffs:
(i) the variational bound is \emph{strictly tighter} than the
fully-labelled ELBO at any non-Bayes-optimal class posterior, because
the analytic marginalisation of $c$ and $(g, e)$ replaces two Jensen
inequalities with equalities;
(ii) the labelled-variant cross-column hard-zero structural
constraint disappears, so per-column factorised $q^{(\tau)}$ is
adequate (the column-Markov-chain upgrade is no longer required);
(iii) per-branch evaluation is roughly
$|\mathcal{T}^{\text{lab}}|^2 / |\mathcal{T}|^2$-times cheaper.

The bound is $\log p(\text{MSA} \mid \parsetree, \theta) \geq
\log \tilde{p}(\text{MSA} \mid \parsetree, \theta) \geq
\mathcal{L}[q]$, in the same restricted-model sense as
appendix~\ref{sec:varanc-presence}; the
restriction-gap ($q$-independent) and variational-gap (KL)
decomposition carries over.

\subsection{Setting and reduced state space}
\label{sec:varanc-mixdom-setting}

Fix a rooted phylogeny $\parsetree$ with internal-node set
$\mathcal{I}$, leaf set $\mathcal{L}$, and branch lengths
$\branchlen_e$ on each edge $e$. The MSA has $L$ columns; each leaf
$v \in \mathcal{L}$ has an observed presence indicator
$X^v_n \in \{0, 1\}$ and (when present) a residue
$\anctok^v_n \in \alphabet$.

The per-(node, column) variational state space is
\begin{equation}
\label{eq:Z-mixdom}
\mathcal{Z} \;=\; \{N, D\} \cup \mathcal{T},
\qquad
\mathcal{T} \;=\; [N_{\text{dom}}] \times [N_{\text{fr}}].
\end{equation}
A tuple-valued state $\tau = (d, f) \in \mathcal{T}$ encodes only the
column's domain $d$ and fragchar $f$; the per-character class $c$ is
marginalised at the model level (Section~\ref{sec:varanc-mixdom-substitution})
and the fragchar-boundary / domain-end indicators $(g, e)$ are
absorbed into the reduced WFST kernel
(Section~\ref{sec:varanc-mixdom-reduced-wfst}).

The presence indicator at internal node $v$, column $n$ is
$X^v_n = \mathbb{1}\{Z^v_n \in \mathcal{T}\}$. Leaf states are
partially clamped: $X^v_n = 0$ pins $Z^v_n \in \{N, D\}$;
$X^v_n = 1$ pins $Z^v_n \in \mathcal{T}$ (consistent with the
column-wide tuple $\tau_n$ — see below).

\subsection{Restricted generative model}
\label{sec:varanc-mixdom-restricted}

Following Section~\ref{sec:varanc-presence-restricted}, the model
joint over MSA columns and internal-node states is
\begin{equation}
\label{eq:joint-mixdom}
p(\text{MSA},\, \{Z^v\}_{v \in \mathcal{I}} \mid \parsetree, \theta)
\;=\; p_{\text{singlet}}^{\text{red}}(Z^{\text{root}})
\,\prod_{(v \to w) \in \parsetree}
   \hat{P}^{\text{WFST}}\!(Z^w \mid Z^v, \branchlen_{vw}, \theta)
\,\prod_{n: \mathcal{F}_n \neq \emptyset}
   L^{(\text{sub}),\text{tot}}_n(f_n; \mathcal{F}_n),
\end{equation}
where $\mathcal{F}_n$ is the Fitch-parsimony subtree of column $n$,
$L^{(\text{sub}),\text{tot}}_n(f; \mathcal{F}_n)$ is the
class-marginalised Felsenstein column-substitution likelihood under
fragchar $f$ (Section~\ref{sec:varanc-mixdom-substitution}), and
$p_{\text{singlet}}^{\text{red}}, \hat{P}^{\text{WFST}}$ are the
\emph{reduced} singlet HMM and WFST defined on $(d, f)$ alone
(Section~\ref{sec:varanc-mixdom-reduced-wfst}). As in the simple
case we marginalise~\eqref{eq:joint-mixdom} over internal patterns
supported on the $L$ observed columns to obtain
$\tilde{p}(\text{MSA} \mid \parsetree, \theta) \le
p(\text{MSA} \mid \parsetree, \theta)$ (with the same ghost-column
caveat), and bound the latter from below.

\subsection{Variational family}
\label{sec:varanc-mixdom-q}

The variational $q$ factorises over MSA columns,
$q(\{Z^v_n\}) = \prod_{n=1}^L q_n$, and within each column,
\begin{equation}
\label{eq:q-mixdom-factor}
q_n \;=\; q_n^{(\tau)}(\tau_n) \cdot
          q_n^{(\pi|\tau)}\!\big(\{Z^v_n\}_{v \in \mathcal{I}} \,\big|\, \tau_n\big),
\end{equation}
where $q_n^{(\tau)}$ is a free categorical over $\mathcal{T}$ on the
column-wide tuple, and $q_n^{(\pi|\tau)}$ is the same 3-state
irreversible tree-structured graphical model of
Section~\ref{sec:varanc-presence-q} (with $\state{P} \equiv \tau_n$).
By construction $q$ assigns zero mass to within-column configurations
in which Present nodes carry differing tuples — the column has one
$(d, f)$, full stop. The cross-column structural-label constraints
that forced a column-Markov chain in the labelled formulation
\emph{vanish} under reduction
(Section~\ref{sec:varanc-mixdom-cross-column-vanishes}).

\paragraph{Free parameter count.}
Per column: $|\mathcal{T}| - 1 = N_{\text{dom}} N_{\text{fr}} - 1$
free parameters in $q_n^{(\tau)}$, plus
$2 |\mathcal{E}|$ in the inner 3-state graphical model
(as in Section~\ref{sec:varanc-presence-special-cases}), plus 2 for
the inner root distribution. For typical
$(N_{\text{dom}}, N_{\text{fr}}) = (3, 2)$ that is 5 tuple
parameters per column, vs $4 N_{\text{cl}} N_{\text{fr}} N_{\text{dom}} - 1$
in the labelled formulation (e.g.\ 239 for $N_{\text{cl}} = 10$) and
$\sim |\mathcal{T}|^2$ in the labelled column-Markov upgrade
($\sim 6 \times 10^4$).

\subsection{\texorpdfstring{Reduced WFST: marginalising $(g, e)$ and the class $c$}{Reduced WFST: marginalising (g, e) and the class c}}
\label{sec:varanc-mixdom-reduced-wfst}

The reduced per-character WFST kernel is
\begin{equation}
\label{eq:reduced-WFST}
\hat{T}_{ss'}\!\big((d, f), (d', f');\, \branchlen, \theta\big),
\qquad s, s' \in \{\sta, \mat, \ins, \del, \fin\},
\end{equation}
obtained by marginalising the labelled-MixDom WFST
$T^{\text{lab}}_{ss'}((c, f, g, d, e), (c', f', g', d', e');
\branchlen, \theta)$ over $(c, c', g, e, g', e')$ at each end. The
reduced state space is $\{\sta, \fin\} \cup \{(s, d, f) :
s \in \{\mat, \ins, \del\}\}$ — i.e.\
$3 N_{\text{dom}} N_{\text{fr}} + 2$ states.

\paragraph{Routes from $(d, f)$ to $(d', f')$.}
Under the labelled MixDom model, the per-character labelled
transition $(d, f) \to (d', f')$ admits up to three latent routes,
indexed by the value of $(g, e)$ at the source position:
\begin{enumerate}[nosep, label=R\arabic*.]
\item \emph{Intra-fragment fragchar transition} ($g = 0$): the
current character is not the last of its fragment, and the next
character stays in the same fragment with fragchar transitioning
$f \to f'$ via $\ext^{(d)}_{f, f'}$. Restricted to $d' = d$.
\item \emph{New fragment, same domain} ($g = 1, e = 0$): the
current fragment terminates ($\notext^{(d)}_f$) and the same domain
starts a new fragment ($\kappa_d$) with first fragchar drawn from
$\fragdist_{d, f'}$. Restricted to $d' = d$.
\item \emph{New domain} ($g = 1, e = 1$): the current fragment
terminates, the current domain ends ($1{-}\kappa_d$), and (after
possibly skipping geometrically-many empty domains) a new domain
$d'$ begins with its first fragment's first fragchar drawn from
$\fragdist_{d', f'}$. Available for any $d'$, including the
self-recurrence $d' = d$.
\end{enumerate}

For $d' \neq d$ only R3 is enabled, so the source's $(g, e) = (1, 1)$
is uniquely determined by the transition. For $d' = d$, multiple
routes contribute and the latent $(g, e)$ at the source carries a
non-trivial posterior over routes given the observed transition.
The diagonal-extension special case ($\ext^{(d)}_{f, f'} = \ext_f
\delta_{f, f'}$) suppresses R1 whenever $f' \neq f$, but R2 and R3
still mix for any $(d, f) \to (d, f')$, so the route-sum machinery
introduced below is required even there. With a general off-diagonal
$\ext^{(d)}$ matrix, R1 admits $f' \neq f$ transitions, so
the route-sum acquires a third contributing term.

The per-route singlet contribution, accounting for the joint prior
on $(g, e)$ at the source and the singlet emission of the next
character, is:
\begin{align}
\label{eq:omega-routes}
\omega^{(R1)}_{(d, f) \to (d', f')}
&= \delta_{d, d'}\,\ext^{(d)}_{f, f'}, \\
\omega^{(R2)}_{(d, f) \to (d', f')}
&= \delta_{d, d'}\,\notext^{(d)}_f\,\kappa_d\,\fragdist_{d, f'}, \\
\omega^{(R3)}_{(d, f) \to (d', f')}
&= \frac{\notext^{(d)}_f\,(1-\kappa_d)\,\kappa_\main\,\domdist_{d'}\,\kappa_{d'}\,\fragdist_{d', f'}}{1-\zeta},
\end{align}
where $\zeta = \kappa_\main \sum_{d''} \domdist_{d''}(1-\kappa_{d''})$
is the empty-domain renormalisation. Each $\omega^{(r)}$
factorises as $P((g, e) = (g_r, e_r) \mid (d, f)) \cdot
P_{\text{singlet}}((d', f') \mid (d, f), g_r, e_r)$, so summing
over routes recovers the marginal singlet emission probability:
\begin{equation}
\label{eq:omega}
\omega^{(d, f, d', f')}
\;=\; \sum_r \omega^{(r)}_{(d, f) \to (d', f')}
\;=\; \delta_{d, d'}\!\Big[\ext^{(d)}_{f, f'}
                          + \notext^{(d)}_f\,\kappa_d\,\fragdist_{d, f'}\Big]
\;+\; \frac{\notext^{(d)}_f\,(1-\kappa_d)\,\kappa_\main\,\domdist_{d'}\,\kappa_{d'}\,\fragdist_{d', f'}}{1-\zeta}.
\end{equation}

\paragraph{Marginalisation of $c$.}
The class $c$ at column $n$ is generated once per column by
$\classdist_{f_n, \cdot}$ at column birth and governs the
substitution likelihood across all branches. The labelled WFST's
indel block is class-independent, so $\sum_c \classdist_{f, c} = 1$
trivially. Class-marginalisation of $\hat{T}$ has no residue at the
indel level; it is handled separately in the substitution term
(Section~\ref{sec:varanc-mixdom-substitution}).

\subsubsection{Step-by-step derivation of the reduced kernel}
\label{sec:varanc-mixdom-derivation}

The route-sum~\eqref{eq:T-hat} is not stipulated; it follows from
two facts about the labelled MixDom model. We state and prove
both.

\paragraph{Notation.} Write $Z_n = (d_n, f_n, g_n, e_n)$ for the
labelled chain state at position $n$ (dropping $c_n$, which is
class-independent for the indel block) and $\tau_n = (d_n, f_n)$
for the reduced state. Let $\sigma$ denote the labelled singlet
HMM transition kernel, and let $T^{\text{lab}}_{ss'}$
denote the labelled conditional WFST kernel of the full labelled-MixDom
model (the per-character transducer on the labelled state
$(c, f, g, d, e)$, before the $(c, g, e)$-marginalisation performed
here). The labelled per-character Pair HMM
joint factorises as
\begin{equation}
\label{eq:pair-fact}
P\big((s', Z') \mid (s, Z)\big)
\;=\; T^{\text{lab}}_{ss'}\!\big(Z, Z'\big)\,
      \sigma\!\big(Z' \mid Z\big),
\end{equation}
i.e.\ the WFST conditional weight times the singlet emission
weight --- the \textit{Singlet} $\circ$ \textit{WFST} $=$ \textit{Pair HMM}
factorisation of the labelled-MixDom model.

\paragraph{Singlet-table convention.}
Throughout the proof we use the labelled singlet HMM table with the
following normalisation. Two structural
facts about the labelled MixDom chain underpin this:
\begin{enumerate}[nosep, label=(\alph*)]
\item \emph{$e$ is fragment-level.} The domain-end indicator $e$
is constant for all characters within a single fragment (it
records whether the current fragment is the last in its domain,
which is set when the fragment begins and does not change as the
fragment extends). Hence within-fragment transitions ($g{=}0$
rows) have $e' = e$ implicitly, and a single ``$g{=}0$ row''
covers all source $e$ values without depending on $e$.
\item \emph{Destination indicators $(g', e')$ are summed implicitly.}
The destination $(g', e')$ in each row is treated as a free index
ranging over $\{0, 1\}^2$, with the implicit prior
$\pi(g', e' \mid d', f')$ to be applied separately if a
finer-grained joint over destination indicators is needed.
\end{enumerate}
Under these conventions, the table entry $\widetilde\sigma$ for a
given source-row $(d, f, g, e)$ and destination structural state
$(d', f')$ is the joint
\[
\widetilde\sigma((d, f, g, e) \to (d', f'))
\;=\; \pi(g, e \mid d, f)\,\cdot\,\sigma_{\text{cond}}((d', f')
\mid (d, f, g, e)),
\]
where $\sigma_{\text{cond}}$ is the source-conditional next-character
emission probability and $\pi(g, e \mid d, f)$ is the source-row
marginal. As a sanity check: summing each row of the table over
$(d', f')$ recovers exactly $\pi(g, e \mid d, f)$, which sums in turn
to~1 over the three rows (with termination column included).
For the $g{=}0$ row: $\sum_{f'} \ext^{(d)}_{f, f'} = 1{-}\notext^{(d)}_f
= \pi(g{=}0 \mid d, f)$. For the $g{=}1, e{=}0$ row:
$\sum_{f'} \notext^{(d)}_f \kappa_d \fragdist_{d, f'} = \notext^{(d)}_f
\kappa_d = \pi(g{=}1, e{=}0 \mid d, f)$. For the $g{=}1, e{=}1$ row,
the row sum (over destinations
plus termination) equals $\notext^{(d)}_f (1{-}\kappa_d)
= \pi(g{=}1, e{=}1 \mid d, f)$.

\paragraph{Lemma 1 (Conditional independence of $(g, e)$).}
\emph{Under the labelled MixDom singlet construction, $(g_n, e_n)
\perp \{Z_k : k < n\} \mid (d_n, f_n)$ for every $n \geq 1$, with}
\begin{align}
\label{eq:ge-prior}
\pi(g{=}0 \mid d, f) &= 1 - \notext^{(d)}_f, \\
\pi(g{=}1, e{=}0 \mid d, f) &= \notext^{(d)}_f \,\kappa_d, \\
\pi(g{=}1, e{=}1 \mid d, f) &= \notext^{(d)}_f\,(1 - \kappa_d).
\end{align}

\emph{Proof.} By the singlet-table convention above, the
labelled-chain joint factorises position-by-position as
\[
P(Z_n \mid Z_{n-1}) = P(d_n, f_n \mid Z_{n-1}) \cdot \pi(g_n, e_n \mid d_n, f_n),
\]
because the table entry at position $n$ has the destination prior
$\pi(g_n, e_n \mid d_n, f_n)$ as an explicit multiplicative factor,
independent of $Z_{n-1}$
beyond what flows through $(d_n, f_n)$. Marginalising the chain
joint over $Z_{1:n-1}$ at fixed $(d_n, f_n)$ leaves the
$\pi(g_n, e_n \mid d_n, f_n)$ factor untouched, establishing
$P(g_n, e_n \mid d_n, f_n, Z_{1:n-1}) = \pi(g_n, e_n \mid d_n, f_n)$.
Substituting the explicit Bernoulli forms
($g_n \sim \text{Bern}(\notext^{(d_n)}_{f_n})$ and, given $g_n = 1$,
$e_n \sim \text{Bern}(1 - \kappa_{d_n})$, with $e_n$ at $g_n = 0$
folded into the joint as the residual~$\notext^{(d)}_f$ row)
gives~\eqref{eq:ge-prior}. \qedhere

\paragraph{Lemma 2 (Singlet route-decomposition).}
\emph{The marginal singlet emission $\omega^{(d, f, d', f')} :=
\sum_{g, e, g', e'} \widetilde\sigma((d, f, g, e) \to (d', f', g', e'))$
decomposes as $\omega = \omega^{(R1)} + \omega^{(R2)} + \omega^{(R3)}$,
with the per-route weights of eq.~\eqref{eq:omega-routes}.}

\emph{Proof.} The labelled singlet table partitions the source
$(g, e)$ values into three rows:
\begin{itemize}[nosep]
\item Row $g{=}0$ (combining all source $e$ values, since the
$g{=}0$ row weight does not depend on source $e$): table entry
$\widetilde\sigma_{R1} = \ext^{(d)}_{f, f'}$ for destination
$(d, f', g', e')$ summed over destination indicators.
\item Row $g{=}1, e{=}0$: $\widetilde\sigma_{R2} = \notext^{(d)}_f
\kappa_d \fragdist_{d, f'}$.
\item Row $g{=}1, e{=}1$: $\widetilde\sigma_{R3} = \notext^{(d)}_f
(1{-}\kappa_d) \kappa_\main \domdist_{d'} \kappa_{d'} \fragdist_{d', f'}
/ (1{-}\zeta)$.
\end{itemize}
By the singlet-table convention, each row is the \emph{joint}
contribution including the source-$(g, e)$ prior; in particular
$\widetilde\sigma_{R1} = \pi(g{=}0 \mid d, f) \cdot
\sigma_{\text{cond}}((d, f', g'{=}\cdot, e'{=}\cdot) \mid (d, f,
g{=}0, e))$ summed over destination $(g', e')$. Thus
$\omega^{(r)} = \widetilde\sigma_r$ directly: the row weight already
equals the route contribution to the marginal singlet emission, so
no separate ``prior cancellation'' step is needed --- the source
prior was baked in from the start. Summing the three rows gives
$\omega^{(d, f, d', f')}$ in~\eqref{eq:omega}. The fact that
each row has support only on destinations consistent with its
route is by inspection of the singlet table. \qedhere

\paragraph{Proposition (Reduced kernel).}
\emph{The reduced per-character marginal Pair HMM kernel,
$\hat T_{ss'}((d, f), (d', f')) := \sum_{g, e, g', e'} \pi(g, e \mid d, f)
\cdot P((s', d', f', g', e') \mid (s, d, f, g, e))$,
satisfies the route-sum~\eqref{eq:T-hat}.}

\emph{Proof.} Substitute the Pair HMM factorisation
\eqref{eq:pair-fact}, and replace
$\pi(g, e \mid d, f) \cdot \sigma_{\text{cond}}(\cdots)
= \widetilde\sigma(\cdots)$ via the singlet-table convention:
\begin{align*}
\hat T_{ss'}((d, f), (d', f'))
&= \sum_{g, e, g', e'} \pi(g, e \mid d, f)\,
   T^{\text{lab}}_{ss'}\!\big((d, f, g, e), (d', f', g', e')\big)\,
   \sigma_{\text{cond}}\!\big((d', f', g', e') \mid (d, f, g, e)\big) \\
&= \sum_{g, e, g', e'}
   T^{\text{lab}}_{ss'}\!\big((d, f, g, e), (d', f', g', e')\big)\,
   \widetilde\sigma((d, f, g, e) \to (d', f', g', e')).
\end{align*}
By Lemma~2, $\widetilde\sigma$ has support partitioned into the
three routes $R1, R2, R3$, each with a unique source $(g_r, e_r)$.
The labelled WFST entry for fixed source $(g_r, e_r)$ and
structural transition $(s, s')$ depends on the destination
$(g', e')$ in only two ways: (i)~the label-preservation constraint
on $\mat$ (which forces $g' = g_{r,\text{dest}}$ and $e' = e_{r,\text{dest}}$
where the destination-side $(g_r, e_r)$ are determined by the route
and the destination structural label, trivialising the destination
sum); and (ii)~for $\ins$, the destination prior on $(g', e')$ at
the new $(d', f')$ is folded into the WFST as a $\pi(g', e' \mid d',
f')$ factor (the WFST, defined as the labelled Pair~HMM joint divided
by the labelled singlet emission, makes this division explicit). In both
cases the $(g', e')$ summation cleanly factors out:
\[
\sum_{g', e'} T^{\text{lab}}_{ss'}((d, f, g_r, e_r),
(d', f', g', e'))\,\pi(g', e' \mid d', f')^{-1}\,\pi(g', e' \mid d',
f') = \tilde T^{\text{lab}, (r)}_{ss'}((d, f), (d', f'))
\]
(the inverse-prior in the WFST cancels the prior in
$\widetilde\sigma$'s destination factor; for $\mat$ the inverse
prior is trivially~1 because of label preservation). Combining
with the source weights $\widetilde\sigma_r$ summed over destinations
giving $\omega^{(r)}$ from Lemma~2, we obtain eq.~\eqref{eq:T-hat}.
\qedhere

\paragraph{Corollary (Markov property of reduced chain).}
\emph{The marginal chain in $(d, f)$ obtained by summing $(g, e)$
out of the labelled singlet chain is exactly Markov, with
per-step kernel $\hat T_{\sta s}$ at the reduced state space.}

\emph{Proof.} We show $P(\tau_n \mid \tau_{1:n-1}) = P(\tau_n \mid
\tau_{n-1})$ for all $n \geq 2$, which is the defining property of
a Markov chain. By the singlet's chain factorisation, $P(\tau_n
\mid Z_{1:n-1}) = \sigma_{\text{cond, marg}}(\tau_n \mid Z_{n-1})$,
which depends on $Z_{n-1} = (\tau_{n-1}, g_{n-1}, e_{n-1})$.
Marginalising $(g_{n-1}, e_{n-1})$ against its conditional
distribution given $\tau_{1:n-1}$:
\[
P(\tau_n \mid \tau_{1:n-1}) = \sum_{g_{n-1}, e_{n-1}}
   P(g_{n-1}, e_{n-1} \mid \tau_{1:n-1})\,
   \sigma_{\text{cond, marg}}(\tau_n \mid \tau_{n-1}, g_{n-1}, e_{n-1}).
\]
By Lemma~1 (Markov property of $(g, e)$), $P(g_{n-1}, e_{n-1} \mid
\tau_{1:n-1}) = \pi(g_{n-1}, e_{n-1} \mid \tau_{n-1})$, depending
only on $\tau_{n-1}$. Substituting:
\[
P(\tau_n \mid \tau_{1:n-1}) = \sum_{g_{n-1}, e_{n-1}}
   \pi(g_{n-1}, e_{n-1} \mid \tau_{n-1})\,
   \sigma_{\text{cond, marg}}(\tau_n \mid \tau_{n-1}, g_{n-1}, e_{n-1})
\]
which is $\hat T_{\sta s}$ (eq.~\eqref{eq:T-hat} for the $s = s' = $
trivial-transition case, i.e.\ depending only on $(\tau_{n-1},
\tau_n)$). This depends only on $\tau_{n-1}$, establishing the
Markov property. The product factorisation $P(\tau_1, \ldots,
\tau_L) = \prod_n \hat T(\tau_{n-1}, \tau_n)$ follows by chain
rule. \qedhere

\paragraph{Caveats.} Two technical conditions buried in the proof
deserve explicit flagging:
\begin{enumerate}[nosep]
\item \emph{Independence of WFST weights from destination $(g', e')$:}
the labelled WFST entries $T^{\text{lab}}_{ss'}((d, f, g, e),
(d', f', g', e'))$ depend on the structural transition type
($\mat$/$\ins$/$\del$ at both ends) and on
the source $(g, e)$ via the case split into mid-fragment /
fragment-boundary / domain-boundary, but they do not depend on
the destination $(g', e')$ except through the trivial label-preservation
($\mat$ requires matching labels) constraint already implicit in
$T^{\text{lab}}$. This is by inspection of the labelled-MixDom WFST
tables.
\item \emph{Boundary entries:} the proof above covers
non-boundary positions $1 \leq n \leq L$ with $(s, s') \in \{\mat,
\ins, \del\}$. The $\sta$ and $\fin$ row/column entries of $\hat T$
require separate derivation because the $\sta$ source has no
preceding $(g, e)$ to enumerate routes over, and the $\fin$
destination requires the singlet's termination weight rather than a
next-character emission. These are flagged as
Section~\ref{sec:varanc-mixdom-open-issues}, item~I.2.
\end{enumerate}

\paragraph{Markov property of the reduced chain (consequence).}
The latent $(g, e)$ at position $n$ depends only on $(d_n, f_n)$
and the model parameters (Lemma~1). Consequently, summing out
$(g, e)$ at each position yields a marginal chain in $(d, f)$ that
remains \emph{exactly Markov}, with per-step transition kernel
$\hat T$ defined below. This justifies the path-LL formulation in
Section~\ref{sec:varanc-mixdom-pathLL} as exact (not a
mean-field approximation) under the reduced state space.

\paragraph{Reduced kernel as a route-sum.}
Combining the route enumeration with the labelled WFST entries
gives the reduced per-character kernel as a sum of route
contributions:
\begin{equation}
\label{eq:T-hat}
\boxed{\;
\hat{T}_{ss'}\!\big((d, f), (d', f');\, \branchlen, \theta\big)
\;=\; \sum_{r \in \mathcal{R}((d, f), (d', f'))}
   \omega^{(r)}_{(d, f) \to (d', f')}\,
   \tilde{T}^{\text{lab}, (r)}_{ss'}\!\big((d, f), (d', f');\, \branchlen, \theta\big),
\;}
\end{equation}
with route set $\mathcal{R}((d, f), (d', f')) = \{R3\}$ for
$d' \neq d$ and $\{R1, R2, R3\}$ for $d' = d$, and per-route
labelled WFST entry
\begin{equation}
\label{eq:tilde-T-lab-r}
\tilde{T}^{\text{lab}, (r)}_{ss'}\!\big((d, f), (d', f');\, \branchlen, \theta\big)
\;=\; \sum_{g', e'} T^{\text{lab}}_{ss'}\!\big((d, f, g_r, e_r), (d', f', g', e');\, \branchlen, \theta\big),
\end{equation}
where $T^{\text{lab}}$ is the \emph{conditional} labelled WFST,
defined as the labelled Pair HMM
joint divided by the labelled singlet emission, and
$(g_r, e_r)$ is the source $(g, e)$ for route $r$:
$(g_{R1}, e_{R1}) = (0, e)$ for any $e$ (since within-fragment
transitions do not constrain $e$; the labelled-singlet table treats
$e$ as a free variable in the
$g{=}0$ row, so we adopt $(0, 0)$ as the canonical representative);
$(g_{R2}, e_{R2}) = (1, 0)$; and $(g_{R3}, e_{R3}) = (1, 1)$.
The product $\omega^{(r)} \cdot \tilde{T}^{\text{lab}, (r)}_{ss'}$
then equals the per-route contribution to the labelled Pair HMM
joint $(d, f) \to (d', f')$ summed over destination indicators
$(g', e')$, so summing over routes recovers the marginal Pair HMM
joint at the reduced state space; division of $T^{\text{lab}}$
by the singlet was applied per-route inside $\omega^{(r)}$ to
avoid double-counting.

\paragraph{When the route-sum collapses.}
The previously-claimed factorisation $\hat T = \omega \cdot
\tilde T^{\text{lab}}$ holds only if, for the given $(s, s')$, the
labelled WFST entry $\tilde T^{\text{lab}, (r)}_{ss'}$ is independent
of $r$ across the enabled routes. This essentially never happens for
the indel block. For $\mat\to\mat$, the
per-route conditional WFST weights (Pair HMM joint divided by the
route's singlet emission) are
\[
W^{(R1)}_{\mat\mat} = 1, \qquad
W^{(R2)}_{\mat\mat} = (1{-}\beta_d)\alpha_d, \qquad
W^{(R3)}_{\mat\mat} \;\propto\; \tau^{(d)}_{\mat,\fin}\,
   \nonemptytrans^{[d \to d']}_{\mat,\mat}\,\tau^{(d')}_{\sta,\mat},
\]
where $W^{(R3)}_{\mat\mat}$ involves the destination-domain-specific
top-level effective transition (in the implementation,
\texttt{T\_exit\_k}; structurally, $\nonemptytrans^{[d\to d']}_{\mat,\mat}$
is the destination-typed version of the empty-domain-summed top-level
$\mat\to\mat$ kernel, with the $\domdist_{d'}$ and other normalisation
factors absorbed).
The three weights are distinct in general:
R1 carries no separate BDI factor because in-fragment extension is
generated by the singlet alone (the descendant character matches
trivially via fragment continuation); R2 carries the standard
new-fragment match factor $(1{-}\beta_d)\alpha_d$ (the singlet's
$\kappa_d$ cancels into $\tau^{(d)}_{\mat,\mat} = (1{-}\beta_d)\alpha_d\kappa_d$);
R3 carries the cross-domain entry weight, involving both the source
domain's exit $\tau^{(d)}_{\mat,\fin} = (1{-}\beta_d)(1{-}\kappa_d)$
and the destination domain's TKF92 entry
$\tau^{(d')}_{\sta,\mat} = (1{-}\beta_{d'})\alpha_{d'}\kappa_{d'}$.
Hence the route-sum genuinely cannot be collapsed to a single
labelled-WFST entry for any $(d, f) \to (d', f')$ transition that
admits more than one route, which includes every same-domain
transition $(d, f) \to (d, f')$.

The reduced kernel is no more row-stochastic than the labelled one
was --- both are conditional transducers --- but the
input-conditional normalisation of the transducer is preserved by
the route-sum marginalisation.

\paragraph{Numerical verification.}
The route-sum
\eqref{eq:T-hat} has been verified at $t = 0.1$ on a
$(N_\text{dom}, N_\text{fr}) = (2, 2)$ instance (random parameters,
fixed seed) by reconstructing every $\mat\!\to\!\mat$ entry of the
Pair HMM as $\sum_r \omega^{(r)}\,W^{(r)}_{\mat\mat}$ and comparing
against the corresponding entry of \texttt{build\_nested\_trans}; the
maximum absolute discrepancy across the 16-entry $\mat\!\to\!\mat$
block is $1.1 \times 10^{-16}$, i.e.\ floating-point noise. The
script is at \texttt{python/verify\_reduced\_wfst\_routes.py}. The
previous $t = 0$ check did not test the inter-route discrepancy in
$W^{(r)}_{ss'}$ because at $t = 0$ the WFST collapses to (near-)identity
and all routes coincide trivially.

\subsection{Per-branch path log-likelihood}
\label{sec:varanc-mixdom-pathLL}

Read parent and child labelled states column-by-column on branch
$v \to w$ to obtain the joint sequence
$(Z^v_n, Z^w_n)_{n=1}^L$. By the same-tuple invariant of $q_n$, only
five joint configurations carry positive variational mass, mapping
to WFST states as in equation~\eqref{eq:state-map} (with $\tau$ now
in the reduced $\mathcal{T}$):
\begin{equation}
\label{eq:state-map-mixdom-reduced}
S_n =
\begin{cases}
\mat & (Z^v_n, Z^w_n) = (\tau, \tau), \\
\del & (Z^v_n, Z^w_n) = (\tau, \state{D}), \\
\ins & (Z^v_n, Z^w_n) = (\state{N}, \tau), \\
\mathsf{Ig} & (Z^v_n, Z^w_n) \in \{(\state{N},\state{N}), (\state{D}, \state{D})\}.
\end{cases}
\end{equation}
With sentinels $S_0 = \sta$, $S_{L+1} = \fin$, the strip-Ignore
reduction of equation~\eqref{eq:branch-LL} carries over:
\begin{equation}
\label{eq:branch-LL-mixdom-reduced}
\log P\!\big(X^w \mid X^v, \branchlen, \theta\big)
\;=\; \sum_{N=1}^{L+1} \delta(S_N \neq \mathsf{Ig})
       \sum_{M=0}^{N-1} \delta(S_M \neq \mathsf{Ig})
         \log \hat{T}_{S_M S_N}\!\big(\tau_{n(M)},\, \tau_{n(N)};\, \branchlen, \theta\big)
         \prod_{K=M+1}^{N-1} \delta(S_K = \mathsf{Ig}),
\end{equation}
with the reduced WFST $\hat{T}$ from
equation~\eqref{eq:T-hat} replacing the labelled $T^{\text{lab}}$.
Tuples are now in $\mathcal{T} = [N_{\text{dom}}] \times [N_{\text{fr}}]$.

\subsection{\texorpdfstring{Per-column expected indel log-likelihood under $q$}{Per-column expected indel log-likelihood under q}}
\label{sec:varanc-mixdom-elbo-indel}

The variational column-state probabilities $P^{v\to w}_q(S_n = s)$
depend only on the presence/absence pattern through the inner
$q_n^{(\pi|\tau)}$ (which is $\tau$-independent in our factorisation),
so the same belief-propagation machinery from the simple case
(equation~\eqref{eq:column-state-probs}) computes them.

The per-branch expected indel log-likelihood factorises as
\begin{equation}
\label{eq:E-branch-LL-mixdom-reduced}
\expect_q\!\big[\log P(X^w \mid X^v, \branchlen, \theta)\big]
= \sum_{s, s'} \sum_{\tau, \tau' \in \mathcal{T}}
   W^{(v\to w)}_{ss', \tau\tau'}\,
   \log \hat{T}_{ss'}\!\big(\tau, \tau';\, \branchlen, \theta\big),
\end{equation}
with the reduced expected transition counts
\begin{equation}
\label{eq:W-pair-mixdom-reduced}
W^{(v\to w)}_{ss', \tau\tau'}
\;=\; \sum_{N=1}^{L+1}\!\sum_{M=0}^{N-1}
   q^{(\tau)}_{n(M)}(\tau)\, q^{(\tau)}_{n(N)}(\tau')\,
   P^{v\to w}_{q,M}(s)\, P^{v\to w}_{q,N}(s')
   \prod_{K=M+1}^{N-1} P^{v\to w}_{q,K}(\mathsf{Ig}),
\end{equation}
where the per-column factorisation of $q^{(\tau)}$ has reduced the
pairwise tuple weight to a product of per-column marginals (no chain
forward-backward needed). The cumulant trick of
Section~\ref{sec:varanc-presence-cumulant} lifts unchanged to the
inner $K$-product, with stable computation in
$O(L \cdot |s|^2 \cdot |\mathcal{T}|^2)$ per branch via the
running-logaddexp prefix on $\log P^{v\to w}_q(\mathsf{Ig})$. For
$(N_{\text{dom}}, N_{\text{fr}}) = (3, 2)$ this gives
$|\mathcal{T}|^2 = 36$ — versus $|\mathcal{T}^{\text{lab}}|^2 = 57600$
in the labelled variant ($\sim 1600\times$ cheaper per branch).

\subsection{Per-column expected substitution log-likelihood}
\label{sec:varanc-mixdom-substitution}

For each column $n$, define the \emph{Fitch subtree} $\mathcal{F}_n
\subseteq \parsetree$ as the smallest connected subtree containing
all leaves $v$ with $X^v_n = 1$. By stationarity and reversibility of
the substitution model, only nodes inside $\mathcal{F}_n$ contribute
to the substitution likelihood; nodes outside are treated as missing
data and absorbed into normalisation. The Fitch subtree is determined
by the leaf data alone, independently of $q$ and $\tau$.

Because $c$ is a per-column model latent (one class per column,
governing all branches at that column), the column-substitution
likelihood under fragchar $f$ marginalises $c$ at the model level:
\begin{equation}
\label{eq:L-sub-total}
L^{(\text{sub}), \text{tot}}_n(f; \mathcal{F}_n)
\;=\; \sum_{c=1}^{N_{\text{cl}}} \classdist_{f, c}\,
       L^{(\text{sub})}_n(c; \mathcal{F}_n),
\end{equation}
with
\begin{equation}
\label{eq:L-sub-mixdom}
L^{(\text{sub})}_n(c; \mathcal{F}_n)
= \sum_{a \in \alphabet} \eqm^{(c)}_a\,\beta_n^{r_n}(a; c)
\end{equation}
the standard Felsenstein up-pass likelihood at column $n$ on
$\mathcal{F}_n$ under class-$c$'s rate matrix $\revsub^{(c)}$ and
stationary $\eqm^{(c)}$ ($r_n$ is the Fitch-determined root of
$\mathcal{F}_n$, $\beta_n^v(a; c)$ the standard Felsenstein
up-message). The expected substitution
log-likelihood under $q$ involves only the fragchar marginal
$q_n^{(f)}(f) = \sum_d q_n^{(\tau)}(d, f)$:
\begin{equation}
\label{eq:E-sub-mixdom-reduced}
\boxed{\;
\expect_{q_n^{(\tau)}}\!\big[\log L^{(\text{sub}), \text{tot}}_n(f_n; \mathcal{F}_n)\big]
\;=\; \sum_{f=1}^{N_{\text{fr}}} q_n^{(f)}(f)\,
      \log\!\Big[\sum_c \classdist_{f, c}\,
                  L^{(\text{sub})}_n(c; \mathcal{F}_n)\Big].
\;}
\end{equation}
The presence factor $q_n^{(\pi|\tau)}$ does not enter — the
substitution likelihood is fully determined by the (data-determined)
Fitch subtree and the (variational-determined) fragchar marginal.

\paragraph{Why $\log\sum_c \classdist L$, not $\sum_c \classdist \log L$.}
This is the proper integration over the per-column class prior. A
labelled-variant counterpart that carried $c$ as a variational latent
would give $\sum_c q^{(c)}_n(c) \log L^{(\text{sub})}_n(c)$; by Jensen's inequality
$\log \sum_c p_c L_c \geq \sum_c p_c \log L_c$ with $p_c =
\classdist_{f, c}$, equality holding only when
$q^{(c)}_n$ is at the Bayes-optimal class posterior
$q^*(c) \propto \classdist_{f, c} L^{(\text{sub})}_n(c)$.
The reduction analytically marginalises $c$ at the place where the
labelled formulation left an inequality, so the reduced ELBO is at
least as tight as the labelled ELBO and strictly tighter for any
$q^{(c)}$ off the Bayes-optimal class posterior.

\paragraph{Numerical implementation.}
$\log L^{(\text{sub}), \text{tot}}_n(f)$ is computed in log-space:
\begin{equation}
\log L^{(\text{sub}), \text{tot}}_n(f)
\;=\; \mathrm{logsumexp}_c\!\big(\log \classdist_{f, c}
                                   + \log L^{(\text{sub})}_n(c; \mathcal{F}_n)\big),
\end{equation}
with the expected log under $q_n^{(f)}$ a simple weighted sum.

\subsection{ELBO}
\label{sec:varanc-mixdom-elbo}

Combining the four contributions:
\begin{equation}
\label{eq:elbo-mixdom-reduced}
\boxed{\;\;
\log p(\text{MSA} \mid \parsetree, \theta)
\;\geq\;
\log \tilde{p}(\text{MSA} \mid \parsetree, \theta)
\;\geq\;
\mathcal{L}[q],
\;\;}
\end{equation}
with
\begin{equation}
\label{eq:elbo-mixdom-reduced-detail}
\mathcal{L}[q]
= \expect_q\!\big[\log p_{\text{singlet}}^{\text{red}}(Z^{\text{root}})\big]
+ \sum_{(v\to w) \in \parsetree}
   \expect_q\!\big[\log \hat{P}^{\text{WFST}}\!(Z^w \mid Z^v, \branchlen_{vw}, \theta)\big]
+ \sum_{n=1}^L \sum_f q_n^{(f)}(f)\,\log L^{(\text{sub}), \text{tot}}_n(f; \mathcal{F}_n)
+ H[q].
\end{equation}
The first two terms use the reduced singlet HMM (defined on
$\mathcal{T} = [N_{\text{dom}}] \times [N_{\text{fr}}]$ with the
singlet kernel $\omega$ from equation~\eqref{eq:omega}) and the
reduced WFST $\hat{T}$ from
equation~\eqref{eq:T-hat}. The third is the marginalised substitution
term~\eqref{eq:E-sub-mixdom-reduced}.

\paragraph{Entropy decomposition.}
\begin{equation}
\label{eq:entropy-mixdom-reduced}
H[q]
= \sum_{n=1}^L H[q_n^{(\tau)}] \;+\; \sum_{n=1}^L H[q_n^{(\pi|\tau)} \mid \text{MSA}],
\end{equation}
with $H[q_n^{(\tau)}]$ the simple categorical entropy
(closed-form, $|\mathcal{T}| - 1$ free parameters per column) and
$H[q_n^{(\pi|\tau)} \mid \text{MSA}]$ the leaf-conditioned entropy of
the simple-case 3-state graphical model
(equations~\eqref{eq:entropy} and~\eqref{eq:cross-entropy-decomp} of
Section~\ref{sec:varanc-presence-elbo-derivation}, with the same
$\log Z_q$ correction).

\paragraph{Bound interpretation.}
Identical to equation~\eqref{eq:gap-decomp}: the bound is on
$\log p$ with gap = $q$-independent restriction-gap (ghost-column
histories) + variational KL. The reduction $(c, g, e)$-collapse
introduces no additional restriction gap; the analytic
marginalisations are exact at the model level, replacing two
labelled-variant Jensen inequalities with equalities and yielding a
strictly tighter bound at the same variational optimum.

\subsection{Cross-column constraint vanishes}
\label{sec:varanc-mixdom-cross-column-vanishes}

The labelled-variant cross-column constraint (Section M.8 of the
labelled draft) forced $q^{(\tau)}$ to be a column-Markov chain to
avoid hard-zero violations of structural rules
($d_{n+1} = d_n$ when $e_n = 0$, etc.). Under the
$(g, e)$-marginalisation those rules disappear: every entry of
$\omega^{(d, f, d', f')}$ is positive (assuming irreducibility of
$\ext^{(d)}$ and positivity of $\domdist, \fragdist, \notext^{(d)},
\kappa_{\main}, \kappa_{\dom}$), so no $(\tau_n, \tau_{n+1})$ pair
is hard-zero in $\hat{T}$ either. Per-column factorised
$q^{(\tau)}(\tau_1, \ldots, \tau_L) = \prod_n q^{(\tau)}_n(\tau_n)$
puts positive mass everywhere on the model's support and
$\mathbb{E}_q[\log p]$ is finite for any positive $q$.

A column-Markov $q^{(\tau)}$ remains an \emph{optional refinement} —
the true posterior on column-tuples is correlated across columns
even after $(g, e)$-marginalisation (the singlet kernel is genuinely
Markov on $(d, f)$ via $\omega$) — but the per-column factorised
$q^{(\tau)}$ is sufficient for $\mathcal{L}[q]$ to be a finite
proper bound, eliminating the $|\mathcal{T}|^2$ tuple-Markov
parameters per column-pair that the labelled variant required.

\subsection{Special cases and recovery}
\label{sec:varanc-mixdom-reduced-special-cases}

\paragraph{Recovery of simple TreeVarAnc.}
Setting $|\mathcal{T}| = 1$ (single domain, single fragchar) collapses
$\hat{T}$ to the GGI-approximation WFST
$\tkfwfst'$ of Section~\ref{sec:tkf92-wfst} and renders the
substitution term constant, recovering Section~\ref{sec:varanc-presence}
verbatim.

\paragraph{Single-class limit.}
Setting $\classdist_{f, c} = \delta_{c, c_0}$ for a fixed class $c_0$
gives $L^{(\text{sub}), \text{tot}}_n(f) = L^{(\text{sub})}_n(c_0;
\mathcal{F}_n)$, and the substitution term reduces to the standard
Felsenstein log-likelihood under class $c_0$.

\paragraph{Single-fragchar-per-domain limit.}
$N_{\text{fr}} = 1$ collapses $\ext^{(d)}$ to a scalar self-loop, the
within-domain fragchar Markov reduces to a geometric extension, and
the reduced model recovers a simpler scalar-extension form.

\paragraph{Fragment-boundary inference caveat.}
Were the labelled $g$ indicator a deterministic function of
consecutive fragchar transitions ($g_n = \delta(f_n \neq f_{n+1})$
within the same domain --- the special case where intra-fragment
fragchar transitions are forbidden), the variational marginal would
identify fragment boundaries directly. Under the reduced formulation
the variational marginal $q_n^{(\tau)}(d, f)$ at each column gives
the ancestral $(d, f)$ directly, but does \emph{not} uniquely identify
``fragment boundaries'' in that restrictive sense: a fragchar
transition $(d, f) \to (d, f')$ with $f' \neq f$ in the inferred
ancestral trajectory may be either a within-fragment Markov move or
a new-fragment start, depending on the latent route. For ancestral
$(d, f)$
inference at observed columns this distinction is irrelevant. For
downstream tasks that need fragment boundaries explicitly, one
option is to augment $q^{(\tau)}$ post-hoc with a
fragment-boundary indicator inferred from the route posterior over
each $(d, f) \to (d, f')$ transition --- but this is a derived
quantity rather than an independent variational latent.

\subsection{Open issues}
\label{sec:varanc-mixdom-open-issues}

Two points warrant explicit verification before relying on the
construction quantitatively (a third was resolved by numerical
verification at $t=0$, see Section~\ref{sec:varanc-mixdom-reduced-wfst}).

\paragraph{(I.1) $\hat{T}$ row-sum identity.}
The reduced WFST $\hat{T}$ inherits the labelled WFST's
input-conditional normalisation by construction
(Section~\ref{sec:varanc-mixdom-reduced-wfst}), but we have not
explicitly verified the row-sum identity. A small numerical check is
recommended: for $(N_{\text{dom}}, N_{\text{fr}}) \in \{(2, 2),
(3, 2)\}$, build $\hat{T}$ from the labelled $T^{\text{lab}}$ and
verify that composition with the reduced singlet HMM on $(d, f)$
yields the same column-marginal Pair HMM probabilities as the
labelled construction.

\paragraph{(I.2) Start/end-row boundary effects.}
Boundary entries $\hat{T}_{\sta s'}(\sta, \tau')$ and
$\hat{T}_{s\fin}(\tau, \fin)$ involve the labelled WFST's start and
end rows, which carry domain-level boundary indicators
($\nonemptytrans_{\sta\cdot}$, $\nonemptytrans_{\cdot\fin}$) that are
not column-internal. These rows may need explicit derivation rather
than mechanical application of equation~\eqref{eq:T-hat}; the
construction in Section~\ref{sec:varanc-mixdom-reduced-wfst} is
correct for interior $(s, s') \in \{\mat, \ins, \del\}^2$ but should
be verified at the boundaries.

These open issues are flagged here for transparency. Their resolution
is a precondition for production use of the reduced ELBO in
parameter learning; for ancestral-tuple inference at fixed $\theta$
the resulting bias (if any) is $q$-independent and so cancels from
the variational optimum.
